% LISTA DE SÍMBOLOS ------------------------------------------------------------
% Conforme NBR 14724. Símbolos matemáticos recorrentes ao longo do trabalho,
% com significado canônico. Variáveis declarativas operacionais (p10, p50, p90)
% são tratadas como símbolos quantis e não como identificadores de coluna.

\begin{simbolos}
    \item[$y$] Valor observado genérico (retorno realizado) usado em definições de perda
    \item[$y_{\text{true}}$] Valor observado do alvo em uma linha do conjunto fora da amostra
    \item[$q$] Nível-alvo genérico de quantil em definições de perda
    \item[$\hat{q}$] Quantil condicional estimado pelo modelo (genérico)
    \item[$\hat{q}_\tau$] Quantil condicional estimado pelo modelo no nível $\tau$
    \item[$\tau$] Nível de quantil (\textit{quantile level}), $\tau \in (0,1)$
    \item[$p10$] Quantil condicional estimado no nível 0{,}10 (limite inferior do intervalo de predição)
    \item[$p50$] Mediana condicional estimada (quantil no nível 0{,}50)
    \item[$p90$] Quantil condicional estimado no nível 0{,}90 (limite superior do intervalo de predição)
    \item[$L_q(y, \hat{q})$] \textit{Pinball loss} do quantil $q$ avaliada no par $(y, \hat{q})$
    \item[$\mathbb{1}[\cdot]$] Função indicadora: vale 1 se a condição entre colchetes for verdadeira e 0 caso contrário
    \item[$N$] Tamanho da amostra fora da amostra usado no cálculo de métricas
    \item[$i$] Índice de iteração sobre linhas da amostra fora da amostra
    \item[$t$] Índice temporal (instante) em séries de perdas pareadas
    \item[$L_t^{(k)}$] Perda pontual do modelo $k$ no instante $t$
    \item[$d_t$] Diferencial de perdas entre dois modelos no instante $t$
    \item[$\bar{d}$] Média amostral do diferencial de perdas $d_t$
    \item[$\widehat{\mathrm{Var}}(\cdot)$] Estimador de variância de longo prazo (correção HAC)
    \item[$\alpha$] Nível de significância estatística
    \item[$R^2$] Coeficiente de determinação fora da amostra (\textit{out-of-sample})
    \item[$h$] Horizonte de previsão (em passos discretos)
    \item[$h+1$] Horizonte de previsão de um passo à frente
    \item[$h+7$] Horizonte de previsão de sete passos à frente
    \item[$h+30$] Horizonte de previsão de trinta passos à frente (suplementar)
    \item[$\Delta$] Operador de diferença entre variantes ou execuções
    \item[$H_0$] Hipótese nula em teste estatístico
    \item[$H_1$, $H_{2a}$, $H_{2b}$] Hipóteses de pesquisa do trabalho (calibração e superioridade contra \textit{baselines})
\end{simbolos}
